% !TeX program = lualatex
% ================================================================
% Polish Tier 3 Vocabulary Version of DevelopingAveragesStarters
%
% Generated by the server using the following settings:
%
%   language            Polish (pol)
%   text direction      left to right
%   font                Noto Serif
%   fallback font       Noto Serif
%   built from          latex/starters/targeted/KS3/statistics/developingAveragesStarters.tex
%   vocabulary key      latex/starters/targeted/KS3/statistics/developingAveragesStarters/L2/pol/developingAveragesStarters_polishKey.tex
%   tier 3 vocabulary   green   RGB 0 128 0
%   English text        black   RGB 0 0 0
%   Polish text         purple  RGB 112 48 160
%   layout macros       version 3
%
% Please don't edit any information in this comment block - the server
% compares it against the current settings and rebuilds this file when
% they differ, so changes here would be overwritten. However, feel free
% to make updates to the LaTeX below if you want to edit it.
% ================================================================

\documentclass{beamer}
% ================================================================
% Copied from the English sheet this was translated from, so that
% anything it sets up for itself still works here
% ================================================================
\usepackage{tikz}
\usetikzlibrary{calc,arrows.meta,patterns}
\usepackage{multicol}
\usepackage{amsmath}

% ================================================================
% Answer-overlay helpers
% ================================================================
\newcommand{\ablank}[1]{%
  \alt<2>{\textcolor{red}{#1}}{\underline{\phantom{#1}}}%
}
\newcommand{\ashow}[1]{\uncover<2->{\textcolor{red}{\small #1}}}
\newcommand{\ashowq}[1]{\alt<2->{\textcolor{red}{\small #1}}{?}}

% Answers drawn straight on to a TikZ diagram, revealed on overlay 2 like \ashow.
% Space is reserved on overlay 1, so the diagram does not move between slides.
\usetikzlibrary{overlay-beamer-styles}
\tikzset{
  ans/.style     = {red, thick, visible on=<2->},
  ansfill/.style = {red!15, visible on=<2->},
  anslab/.style  = {red, font=\tiny, visible on=<2->},
}

\title{Averages: Mean \& Median}
\author{}
\date{}

% Lego tower: \legotower{x-offset}{height}{colour}
% Each brick is 0.65 wide x 0.55 tall, with two studs
\newcommand{\legotower}[3]{%
  \foreach \i in {1,...,#2}{%
    \draw[fill=#3!55, draw=#3!80!black, line width=0.35pt]
      (#1, {(\i-1)*0.55}) rectangle (#1+0.65, {\i*0.55});
    \draw[fill=#3!80, draw=#3!80!black, line width=0.2pt]
      (#1+0.15,{(\i-1)*0.55+0.38}) circle (0.08);
    \draw[fill=#3!80, draw=#3!80!black, line width=0.2pt]
      (#1+0.44,{(\i-1)*0.55+0.38}) circle (0.08);
  }%
}

% Characters this language's font has not got are borrowed from another,
% rather than being silently dropped - see docs/translated-sheet-latex.md
\directlua{%
  if luaotfload and luaotfload.add_fallback then
    luaotfload.add_fallback("eallatinfallback",
      { "Noto Serif:mode=harf;" })
  else
    tex.error("this luaotfload is too old to register a fallback font")
  end
}
\usepackage[english]{babel}
\babelprovide[import]{polish}
\babelfont[polish]{rm}[Renderer=Harfbuzz, BoldFont={Noto Serif}, ItalicFont={Noto Serif}, BoldItalicFont={Noto Serif}, RawFeature={fallback=eallatinfallback}]{Noto Serif}
\babelfont[polish]{sf}[Renderer=Harfbuzz, BoldFont={Noto Serif}, ItalicFont={Noto Serif}, BoldItalicFont={Noto Serif}, RawFeature={fallback=eallatinfallback}]{Noto Serif}
\newcommand{\ealtext}[1]{\foreignlanguage{polish}{#1}}
\newcommand{\ealtextblock}[1]{{\raggedright\foreignlanguage{polish}{#1}\par}}
% ================================================================
% EAL layout helpers
% ================================================================
\usepackage{xcolor}

\definecolor{ealkeycolour}{RGB}{0,128,0}
\definecolor{ealenglish}{RGB}{0,0,0}
\definecolor{ealtwocolour}{RGB}{112,48,160}

\newsavebox{\ealboxen}
\newsavebox{\ealboxtr}
\newlength{\ealglosswd}
\newlength{\ealglossraise}
\setlength{\ealglossraise}{1.15em}

% a tier 3 word, in English and in translation alike
\newcommand{\ealkey}[1]{\textcolor{ealkeycolour}{#1}}
\newcommand{\ealkeytr}[1]{\textcolor{ealkeycolour}{\ealtext{#1}}}

% One translated word above one English word. Built from kernel
% primitives, never a tabular, which would break wherever the sheet
% loads array, booktabs or tabularx. The box is as wide as the wider of
% the two, so a long translation pushes its neighbours aside instead of
% overlapping them, and it claims its full height so a table row or a
% TikZ node makes room for it.
\newcommand{\ealgloss}[2]{%
  \sbox{\ealboxen}{\ealkey{#1}}%
  \sbox{\ealboxtr}{\small\textcolor{ealtwocolour}{\ealtext{#2}}}%
  \setlength{\ealglosswd}{\wd\ealboxen}%
  \ifdim\wd\ealboxtr>\ealglosswd\setlength{\ealglosswd}{\wd\ealboxtr}\fi
  \leavevmode
  \makebox[\ealglosswd][c]{%
    \makebox[0pt][c]{\usebox{\ealboxen}}%
    \makebox[0pt][c]{\raisebox{\ealglossraise}{\usebox{\ealboxtr}}}%
  }%
}

% opens up line spacing around a block of glosses, so the raised
% translations sit clear of the line above
\newenvironment{ealglossed}{\par\linespread{2.1}\selectfont\ignorespaces}{\par}

% a whole translated sentence above its English counterpart. block level,
% so it does not belong inside a tabular cell
\newcommand{\ealpara}[2]{%
  \par\addvspace{0.5\baselineskip}%
  {\color{ealtwocolour}\ealtextblock{#1}}%
  \penalty200
  {\color{ealenglish}#2\par}%
  \addvspace{0.5\baselineskip}%
}

\begin{document}

\providecommand{\ashow}[1]{\uncover<2->{\textcolor{red}{\small #1}}}
\providecommand{\ashowq}[1]{\alt<2->{\textcolor{red}{\small #1}}{?}}

\begin{frame}{Starter}
\vspace{-0.8em}

\footnotesize
\begin{enumerate}\itemsep2pt
\begin{multicols}{2}
  \item $6 + 9 = \ashowq{15}$
  \item $24 \div 4 = \ashowq{6}$
  \item $7 + 13 + 5 = \ashowq{25}$
  \item $30 \div 5 = \ashowq{6}$
  \end{multicols}
  \item \begin{ealglossed}How do we work out the \ealgloss{mean}{średnia}? \ashow{Add the \ealgloss{values}{wartości} and divide by how many there are.}\end{ealglossed}
  \item \begin{ealglossed}How do we work out the \ealgloss{median}{mediana}? \ashow{Order the \ealgloss{data}{dane} and find the middle \ealgloss{value}{wartość}.}\end{ealglossed}
  \item \begin{ealglossed}Find the \ealgloss{mean}{średnia} of: $3,\ 7,\ 5,\ 9$ \ashow{$=6$}\end{ealglossed}
  \item \begin{ealglossed}Find the \ealgloss{median}{mediana} of: $4,\ 9,\ 2,\ 7,\ 8$ \ashow{$=7$}\end{ealglossed}
  \item \begin{ealglossed}Find the \ealgloss{mean}{średnia} of: $\frac{1}{2},\ \frac{3}{2},\ 2,\ 3$ \ashow{$=\frac{7}{4}$}\end{ealglossed}
  \item \begin{ealglossed}Find the \ealgloss{median}{mediana} of: $-4,\ 2,\ -1,\ -7,\ 1$ \ashow{$=-1$}\end{ealglossed}
  \item \begin{ealglossed}The \ealgloss{mean}{średnia} of $x,\ x+4,\ x+8,\ x+12$ is $15$, what is the total of all the \ealgloss{values}{wartości}? \ashow{$60$}\end{ealglossed}
  \item \begin{ealglossed}Find the \ealgloss{value}{wartość} of $x$ in the previous question \ashow{$x=9$}\end{ealglossed}
\end{enumerate}

\end{frame}


\begin{frame}{Starter -- Solutions}
\vspace{-0.8em}

\footnotesize
\begin{enumerate}\itemsep2pt
\begin{multicols}{2}
  \item $6 + 9 = \mathord{?}$ \textcolor{red}{$15$}
  \item $24 \div 4 = \mathord{?}$ \textcolor{red}{$6$}
  \item $7 + 13 + 5 = \mathord{?}$ \textcolor{red}{$25$}
  \item $30 \div 5 = \mathord{?}$ \textcolor{red}{$6$}
\end{multicols}
  \item \begin{ealglossed}How do we work out the \ealgloss{mean}{średnia}? \textcolor{red}{Add the \ealgloss{values}{wartości} and divide by how many there are}\end{ealglossed}
  \item \begin{ealglossed}How do we work out the \ealgloss{median}{mediana}? \textcolor{red}{Order the \ealgloss{data}{dane} and find the middle \ealgloss{value}{wartość}}\end{ealglossed}
  \item \begin{ealglossed}Find the \ealgloss{mean}{średnia} of: $3,\ 7,\ 5,\ 9$ \textcolor{red}{$=6$}\end{ealglossed}
  \item \begin{ealglossed}Find the \ealgloss{median}{mediana} of: $4,\ 9,\ 2,\ 7,\ 8$ \textcolor{red}{$=7$}\end{ealglossed}
  \item \begin{ealglossed}Find the \ealgloss{mean}{średnia} of: $\frac{1}{2},\ \frac{3}{2},\ 2,\ 3$ \textcolor{red}{$=\frac{7}{4}$}\end{ealglossed}
  \item \begin{ealglossed}Find the \ealgloss{median}{mediana} of: $-4,\ 2,\ -1,\ -7,\ 1$ \textcolor{red}{$=-1$}\end{ealglossed}
  \item \begin{ealglossed}The \ealgloss{mean}{średnia} of $x,\ x+4,\ x+8,\ x+12$ is $15$, what is the total? \textcolor{red}{$60$}\end{ealglossed}
  \item \begin{ealglossed}Find the \ealgloss{value}{wartość} of $x$ \textcolor{red}{$x=9$}\end{ealglossed}
\end{enumerate}
\end{frame}


\begin{frame}{Starter}
\vspace{-0.8em}

\footnotesize
\begin{enumerate}\itemsep2pt
\begin{multicols}{2}
  \item $3 + 7 + 10 = \ashowq{20}$
  \item $48 \div 6 = \ashowq{8}$
  \item $\frac{1}{2} + \frac{1}{4} = \ashowq{\frac{3}{4}}$
  \item $28 \div 4 = \ashowq{7}$
  \end{multicols}
  \item \begin{ealglossed}Which \ealgloss{average}{średnia} requires you to order the \ealgloss{data}{dane} first? \ashow{\ealgloss{Median}{mediana}}\end{ealglossed}
  \item \begin{ealglossed}Find the \ealgloss{median}{mediana} of: $7,\ 2,\ 9,\ 4,\ 5$ \ashow{$=5$}\end{ealglossed}
  \item \begin{ealglossed}Find the \ealgloss{median}{mediana} of: $3,\ 8,\ 1,\ 6,\ 4,\ 10$ \ashow{$=5$}\end{ealglossed}
  \item \begin{ealglossed}Find the \ealgloss{mean}{średnia} and \ealgloss{median}{mediana} of: $2,\ 5,\ 3,\ 8,\ 7$.  
        Which is larger? \ashow{\ealgloss{Mean}{średnia} $=5$, \ealgloss{median}{mediana} $=5$ (equal).}\end{ealglossed}
  \item \begin{ealglossed}Find the \ealgloss{mean}{średnia} and \ealgloss{median}{mediana} of: $-5,\ 1,\ -2,\ 4,\ 7$ \ashow{\ealgloss{Mean}{średnia} $=1$, \ealgloss{median}{mediana} $=1$.}\end{ealglossed}
  \item \begin{ealglossed}The \ealgloss{mean}{średnia} of $a,\ a+3,\ a+6$ is $6$.\end{ealglossed}
        \begin{itemize}\footnotesize
        \item \begin{ealglossed}What is the \ealgloss{median}{mediana} in terms of $a$? \ashow{$a+3$}\end{ealglossed}
          \item \begin{ealglossed}What is the total of the \ealgloss{data}{dane} (use the \ealgloss{mean}{średnia}) \ashow{$18$}\end{ealglossed}
          \item Find $a$ \ashow{$a=3$}
        \end{itemize}
\end{enumerate}

\end{frame}


\begin{frame}{Starter -- Solutions}
\vspace{-0.8em}

\footnotesize
\begin{enumerate}\itemsep2pt
\begin{multicols}{2}
  \item $3 + 7 + 10 = \mathord{?}$ \textcolor{red}{$20$}
  \item $48 \div 6 = \mathord{?}$ \textcolor{red}{$8$}
  \item $\frac{1}{2} + \frac{1}{4} = \mathord{?}$ \textcolor{red}{$\frac{3}{4}$}
  \item $28 \div 4 = \mathord{?}$ \textcolor{red}{$7$}
\end{multicols}
  \item \begin{ealglossed}Which \ealgloss{average}{średnia} requires ordering first? \textcolor{red}{\ealgloss{Median}{mediana}}\end{ealglossed}
  \item \begin{ealglossed}\ealgloss{Median}{mediana} of $7,\ 2,\ 9,\ 4,\ 5$ \textcolor{red}{$=5$}\end{ealglossed}
  \item \begin{ealglossed}\ealgloss{Median}{mediana} of $3,\ 8,\ 1,\ 6,\ 4,\ 10$ \textcolor{red}{$=5$}\end{ealglossed}
  \item \begin{ealglossed}\ealgloss{Mean}{średnia} and \ealgloss{median}{mediana} of $2,\ 5,\ 3,\ 8,\ 7$ \textcolor{red}{\ealgloss{Mean}{średnia} $=5,\ \text{median }=5$}\end{ealglossed}
  \item \begin{ealglossed}\ealgloss{Mean}{średnia} and \ealgloss{median}{mediana} of $-5,\ 1,\ -2,\ 4,\ 7$ \textcolor{red}{\ealgloss{Mean}{średnia} $=1,\ \text{median }=1$}\end{ealglossed}
  \item \begin{ealglossed}\ealgloss{Mean}{średnia} of $a,\ a+3,\ a+6$ is $6$\end{ealglossed}
        \begin{itemize}\footnotesize
          \item \begin{ealglossed}\ealgloss{Median}{mediana} \textcolor{red}{$=a+3$}\end{ealglossed}
          \item Total \textcolor{red}{$=18$}
          \item $a=$ \textcolor{red}{$3$}
        \end{itemize}
\end{enumerate}
\end{frame}


\begin{frame}{Starter}
\vspace{-0.8em}

\footnotesize
\begin{enumerate}\itemsep2pt
\begin{multicols}{2}
  \item $4 + 11 = \ashowq{15}$
  \item $40 \div 8 = \ashowq{5}$
  \item $\frac{3}{4} - \frac{1}{4} = \ashowq{\frac{1}{2}}$
  \item $72 \div 9 = \ashowq{8}$
  \end{multicols}
  \item \begin{ealglossed}A shoe shop owner says,  
        ``The most popular size last week was size 6.''  
    Which \ealgloss{average}{średnia} is she using? \ashow{\ealgloss{Mode}{dominanta}}\end{ealglossed}
  \item \begin{ealglossed}Find the \ealgloss{mode}{dominanta} of: $3,\ 7,\ 3,\ 9,\ 5,\ 3,\ 7$ \ashow{$3$}\end{ealglossed}
  \item \begin{ealglossed}Find the \ealgloss{mode}{dominanta} of: $2,\ 5,\ 8,\ 5,\ 3,\ 8,\ 1$ \ashow{$5$ and $8$}\end{ealglossed}
  \item \begin{ealglossed}Find the \ealgloss{mean}{średnia} and \ealgloss{median}{mediana} of: $4,\ 7,\ 2,\ 9,\ 3$ \ashow{\ealgloss{Mean}{średnia} $=5$, \ealgloss{median}{mediana} $=4$.}\end{ealglossed}
  \item \begin{ealglossed}Find the \ealgloss{mean}{średnia}, \ealgloss{median}{mediana} and \ealgloss{mode}{dominanta} of: $5,\ 3,\ 7,\ 3,\ 5,\ 5,\ 9$ \ashow{\ealgloss{Mean}{średnia} $=\frac{37}{7}$, \ealgloss{median}{mediana} $=5$, \ealgloss{mode}{dominanta} $=5$.}\end{ealglossed}
  \item \begin{ealglossed}The \ealgloss{data}{dane} \ealgloss{set}{zbiór} is $\{4,\ 7,\ 4,\ x,\ 7,\ 3\}$  
        The \textbf{only} \ealgloss{mode}{dominanta} is $4$.  
        Find $x$ and explain your reasoning. \ashow{$x=4$; then $4$ occurs three times.}\end{ealglossed}
\end{enumerate}

\end{frame}


\begin{frame}{Starter -- Solutions}
\vspace{-0.8em}

\footnotesize
\begin{enumerate}\itemsep2pt
\begin{multicols}{2}
  \item $4 + 11 = \mathord{?}$ \textcolor{red}{$15$}
  \item $40 \div 8 = \mathord{?}$ \textcolor{red}{$5$}
  \item $\frac{3}{4} - \frac{1}{4} = \mathord{?}$ \textcolor{red}{$\frac{1}{2}$}
  \item $72 \div 9 = \mathord{?}$ \textcolor{red}{$8$}
\end{multicols}
  \item \begin{ealglossed}Which \ealgloss{average}{średnia} is used? \textcolor{red}{\ealgloss{Mode}{dominanta}}\end{ealglossed}
  \item \begin{ealglossed}\ealgloss{Mode}{dominanta} of $3,\ 7,\ 3,\ 9,\ 5,\ 3,\ 7$ \textcolor{red}{$3$}\end{ealglossed}
  \item \begin{ealglossed}\ealgloss{Mode}{dominanta} of $2,\ 5,\ 8,\ 5,\ 3,\ 8,\ 1$ \textcolor{red}{$5\ \text{and}\ 8$}\end{ealglossed}
  \item \begin{ealglossed}\ealgloss{Mean}{średnia} and \ealgloss{median}{mediana} of $4,\ 7,\ 2,\ 9,\ 3$ \textcolor{red}{\ealgloss{Mean}{średnia} $=5,\ \text{median }=4$}\end{ealglossed}
  \item \begin{ealglossed}\ealgloss{Mean}{średnia}, \ealgloss{median}{mediana} and \ealgloss{mode}{dominanta} of $5,\ 3,\ 7,\ 3,\ 5,\ 5,\ 9$ \textcolor{red}{\ealgloss{Mean}{średnia} $=\frac{37}{7}$, \ealgloss{median}{mediana} $=5$, \ealgloss{mode}{dominanta} $=5$}\end{ealglossed}
  \item \begin{ealglossed}Only \ealgloss{mode}{dominanta} is $4$  
        \textcolor{red}{$x=4$}\end{ealglossed}
\end{enumerate}
\end{frame}


\begin{frame}{Starter}
\vspace{-0.8em}

\footnotesize
\begin{enumerate}\itemsep2pt
\begin{multicols}{2}
  \item $3+4 = \ashowq{7}$
  \item $4 \times 8 = \ashowq{32}$
  \item $-3 + 8 + (-5) = \ashowq{0}$
  \item $\frac{2}{3} + \frac{1}{6} = \ashowq{\frac{5}{6}}$
  \item $-12 \div 4 = \ashowq{-3}$
  \item $\frac{3}{4} \times 8 = \ashowq{6}$
  \end{multicols}
  \item \begin{ealglossed}A class votes for their favourite colour.  
        Can you find the \ealgloss{mean}{średnia}?  
        Which \ealgloss{average}{średnia} should you use instead, and why?  
        \ashow{No; use the \ealgloss{mode}{dominanta}, because the \ealgloss{data}{dane} are categories.}\end{ealglossed}
  \item \begin{ealglossed}Find the \ealgloss{mean}{średnia}, \ealgloss{median}{mediana} and \ealgloss{mode}{dominanta} of:  
        $-3,\ 1,\ -3,\ 4,\ -1,\ 2,\ -3$ \ashow{\ealgloss{Mean}{średnia} $=-\frac{3}{7}$, \ealgloss{median}{mediana} $=-1$, \ealgloss{mode}{dominanta} $=-3$.}\end{ealglossed}
  \item \begin{ealglossed}Find the \ealgloss{median}{mediana} of:  
        $\frac{1}{4},\ \frac{3}{4},\ \frac{1}{2},\ \frac{1}{8},\ \frac{5}{8}$ \ashow{$\frac{1}{2}$}\end{ealglossed}
  \item \begin{ealglossed}The ages in a youth club are:  
        $11,\ 12,\ 12,\ 13,\ 13,\ 13,\ 14,\ 42$
        Which \ealgloss{average}{średnia} would you not use for this example, and why? \ashow{\ealgloss{Mean}{średnia}, because $42$ is an \ealgloss{outlier}{wartość odstająca}.}\end{ealglossed}
  \item \begin{ealglossed}The \ealgloss{mean}{średnia} of $\{x,\ 3,\ x,\ 7,\ 5\}$ is $6$.  
        Find $x$ and state the \ealgloss{mode}{dominanta}. \ashow{$x=6$, \ealgloss{mode}{dominanta} $=6$}\end{ealglossed}
  \item \begin{ealglossed}Five numbers have \ealgloss{mean}{średnia} $8$, \ealgloss{median}{mediana} $7$ and \ealgloss{unique}{jedyna} \ealgloss{mode}{dominanta} of $5$.  
        Two numbers are $5$ and $5$.  
        Find a possible \ealgloss{set}{zbiór}. \ashow{$\{5,5,6,7,17\}$}\end{ealglossed}
\end{enumerate}

\end{frame}


\begin{frame}{Starter -- Solutions}
\vspace{-0.8em}

\footnotesize
\begin{enumerate}\itemsep2pt
\begin{multicols}{2}
  \item $3+4 = \mathord{?}$ \textcolor{red}{$7$}
  \item $4 \times 8 = \mathord{?}$ \textcolor{red}{$32$}
  \item $-3 + 8 + (-5) = \mathord{?}$ \textcolor{red}{$0$}
  \item $\frac{2}{3} + \frac{1}{6} = \mathord{?}$ \textcolor{red}{$\frac{5}{6}$}
  \item $-12 \div 4 = \mathord{?}$ \textcolor{red}{$-3$}
  \item $\frac{3}{4} \times 8 = \mathord{?}$ \textcolor{red}{$6$}
\end{multicols}
  \item \begin{ealglossed}Can you find the \ealgloss{mean}{średnia}?  
        \textcolor{red}{No — \ealgloss{data}{dane} is categorical}\end{ealglossed}
  \item \begin{ealglossed}\ealgloss{Mean}{średnia}, \ealgloss{median}{mediana} and \ealgloss{mode}{dominanta}  
        \textcolor{red}{\ealgloss{Mean}{średnia} $=-\frac{3}{7}$, \ealgloss{median}{mediana} $=-1$, \ealgloss{mode}{dominanta} $=-3$}\end{ealglossed}
  \item \begin{ealglossed}\ealgloss{Median}{mediana}  
        \textcolor{red}{$=\frac{1}{2}$}\end{ealglossed}
  \item \begin{ealglossed}Which \ealgloss{average}{średnia} not to use and why?  
        \textcolor{red}{\ealgloss{Mean}{średnia}, because 42 is an \ealgloss{outlier}{wartość odstająca}}\end{ealglossed}
  \item \begin{ealglossed}$x=?$ and \ealgloss{mode}{dominanta}  
        \textcolor{red}{$x=6$, \ealgloss{mode}{dominanta} $=6$}\end{ealglossed}
  \item \begin{ealglossed}Possible \ealgloss{set}{zbiór}  
        \textcolor{red}{$\{5,5,6,7,17\}$}\end{ealglossed}
\end{enumerate}
\end{frame}


\begin{frame}{Starter}
\vspace{-0.8em}

\footnotesize
\begin{enumerate}\itemsep2pt
\begin{multicols}{2}
\item $3+8 = \ashowq{11}$
\item $120 \div 10 = \ashowq{12}$
  \item $-6 + (-2) + 5 = \ashowq{-3}$
  \item Order from smallest to largest:  
        $\frac{1}{3},\ \frac{5}{6},\ \frac{1}{2},\ \frac{2}{3}$ \ashow{$\frac{1}{3},\ \frac{1}{2},\ \frac{2}{3},\ \frac{5}{6}$}
  \item $3\frac{1}{2} + 2\frac{3}{4} = \ashowq{6\frac{1}{4}}$
  \item $-20 \div 4 = \ashowq{-5}$
  \end{multicols}

  \item \begin{ealglossed}Find the \ealgloss{mean}{średnia}, \ealgloss{median}{mediana} and \ealgloss{mode}{dominanta} of:  
        $-6,\ -2,\ 4,\ -3,\ 7,\ 0,\ -2$ \ashow{\ealgloss{Mean}{średnia} $=-\frac{2}{7}$, \ealgloss{median}{mediana} $=-2$, \ealgloss{mode}{dominanta} $=-2$.}\end{ealglossed}

  \item \begin{ealglossed}True or False?  
        ``Adding 10 to every \ealgloss{value}{wartość} increases the \ealgloss{mean}{średnia} by 10 but leaves the \ealgloss{median}{mediana} unchanged.''  
        \ashow{True}\end{ealglossed}

  \item \begin{ealglossed}The seven \ealgloss{values}{wartości}  
        $n-9,\ n-6,\ n-3,\ n,\ n+3,\ n+6,\ n+9$  
        have a \ealgloss{mean}{średnia} of $4$.  
        State the \ealgloss{median}{mediana} \textbf{without calculation}. \ashow{\ealgloss{Median}{mediana} $=4$}\end{ealglossed}
  \item \begin{ealglossed}A \ealgloss{set}{zbiór} of five \ealgloss{positive}{liczba dodatnia} \ealgloss{integers}{liczby całkowite} has  
        \ealgloss{mean}{średnia} = \ealgloss{median}{mediana} = \ealgloss{mode}{dominanta} = $6$.\end{ealglossed}
        \begin{itemize}\footnotesize
          \item \begin{ealglossed}What is the smallest possible \ealgloss{range}{rozstęp}? \ashow{$1$}\end{ealglossed}
          \item \begin{ealglossed}What is the largest possible \ealgloss{range}{rozstęp}? \ashow{$12$}\end{ealglossed}
        \end{itemize}
\end{enumerate}

\end{frame}


\begin{frame}{Starter -- Solutions}
\vspace{-0.8em}

\footnotesize
\begin{enumerate}\itemsep2pt
\begin{multicols}{2}
  \item $3+8 = \mathord{?}$ \textcolor{red}{$11$}
  \item $120 \div 10$ \textcolor{red}{$12$}
  \item $-6 + (-2) + 5$ \textcolor{red}{$-3$}
  \item Order  
        \textcolor{red}{$\frac13,\ \frac12,\ \frac23,\ \frac56$}
  \item $3\frac12 + 2\frac34$ \textcolor{red}{$6\frac14$}
  \item $-20 \div 4$ \textcolor{red}{$-5$}
\end{multicols}
  \item \begin{ealglossed}\ealgloss{Mean}{średnia}, \ealgloss{median}{mediana}, \ealgloss{mode}{dominanta}  
        \textcolor{red}{\ealgloss{Mean}{średnia} $=-\frac{2}{7}$, \ealgloss{median}{mediana} $=-2$, \ealgloss{mode}{dominanta} $=-2$}\end{ealglossed}
  \item True or False?  
        \textcolor{red}{True}
  \item \begin{ealglossed}\ealgloss{Median}{mediana}  
        \textcolor{red}{$4$}\end{ealglossed}
  \item \begin{ealglossed}\ealgloss{Range}{rozstęp}  
        \textcolor{red}{Smallest $=1$, largest $=12$}\end{ealglossed}
\end{enumerate}
\end{frame}

\end{document}